\section{Introduction}
Web applications remain the primary interface for commerce, identity, and data access; Web Application Firewalls (WAFs) are widely deployed to block common injection classes using signatures and, increasingly, machine-learned components \cite{owasp-crs,waf-survey}. From a defender's standpoint, the decisive question is not merely whether a WAF \emph{exists}, but whether its concrete rules, parsers, and logging configuration provide sufficient coverage against the full range of adversarial inputs a real attacker might send. Signatures can be evaded by semantically equivalent rewrites; edge and application parsers can disagree on decoding, normalization, and parameter binding \cite{mutation-sqli-evasion,parsing-discrepancies,request-smuggling}. Because both phenomena are well documented in the literature, the remaining scientific challenge for defenders is often \emph{measurement}: under what conditions, and for which mutation \emph{shapes}, does WAF policy fail to provide coverage, and what does that mean for rule prioritization?

This paper contributes a \emph{defense-oriented} evaluation that is deliberately operational. We systematically mutate SQLi fragments, probe a production-style AWS WAF deployment in a black-box manner, and analyze coverage gaps at aggregate and variant-family granularity. The implementation is released as a reproducible artifact (\texttt{paper/experiment}), including generators, probe drivers, JSONL traces, and export scripts used to build the tables in Section~\ref{sec:experiments}. We do not propose a new attack tool; we instead supply a \emph{diagnostic methodology} and empirical coverage distributions that security teams can use for targeted rule improvement, blue-team regression testing after policy changes, and benchmarking of future defensive mitigations.

\subsection{Problem Statement}
SQLi evasion exploits at least two mechanisms: (i) \emph{representation-level} mutation that preserves attacker intent while altering substrings, token boundaries, or encodings seen by a pattern matcher; and (ii) \emph{interpretation-level} divergence, where the same octet sequence is decoded or framed differently by the WAF path and the application \cite{akhavani2025waffled}. Defenders must address both to achieve robust coverage, yet current benchmarks rarely supply the family-level granularity needed to know which mechanism is responsible for a given policy failure.

\subsubsection{Mutation-Aware Coverage Gaps}
Prior work treats WAFs as targets for semantic-preserving search. SSQLi frames black-box SQLi as reinforcement learning and reports high bypass rates with SAC \cite{guan2023ssqli}. GSQLi uses GAN-based generation \cite{khan2024gsqli}; WAF-A-MoLE demonstrates mutational fuzzing against ML WAFs \cite{demetrio2020wafamole}. These systems share a common lesson for defenders: policy failure is rarely uniform across operator families---some mutation classes collapse under normalization while others consistently slip through. Understanding which families expose coverage gaps is prerequisite to closing them. Yet published headline rates seldom decompose by operator family on a live managed WAF, leaving defenders without the actionable granularity needed to prioritize rule improvements.

\subsubsection{Discrepancy-Aware Defense Surfaces}
WAFFLED shows that mutating ostensibly benign HTTP structure can induce large sets of parsing-discrepancy bypasses across major WAFs \cite{akhavani2025waffled}. Our live audit narrows scope to URL-visible injection (path or query) with deterministic seed remapping. That choice minimizes one class of discrepancy (parameter-name mismatch) so that residual coverage gaps more directly reflect how the WAF treats alternate encodings and spellings of the same logical injection slot---a controlled setting for studying representation-level policy holes against a fixed rule snapshot.

\subsection{Research Questions}
We organize the empirical work around three diagnostic questions:
\begin{itemize}
  \item \textbf{RQ1 (Coverage baseline).} Under an explicit, script-defined notion of policy gap, what fraction of probed attempts expose a coverage failure in the live AWS WAF lab configuration, and how does that fraction relate to unique payload cardinality?
  \item \textbf{RQ2 (Gap structure).} How much of the coverage-gap mass is concentrated in a small set of mutation families, and which families are effectively defended (e.g., encoding-heavy spawns)? This decomposition directly guides rule prioritization for defenders.
  \item \textbf{RQ3 (Normalization sufficiency).} Does a minimal RFC-oriented percent-decoding fixpoint materially improve a surrogate detector's ability to separate attack-like from benign-like strings---i.e., would trivial decoder alignment alone be sufficient to close the observed coverage gaps?
\end{itemize}

\subsection{Objectives and Contributions}
\begin{enumerate}
\item \textbf{Audit Methodology.} We formalize a black-box diagnostic pipeline---staged variant generation (dialect seeds, AWS-oriented augmentations, learned transforms, optional second-layer transforms), HTTP probing with JSONL logging, and export-driven aggregation---for systematic WAF coverage assessment and policy gap identification.
\item \textbf{Empirical coverage characterization.} On the archived AWS WAF lab run summarized in \texttt{aws\_probe\_summary.json} ($N{=}5{,}000$ variants tested; $N_{\mathrm{gap}}{=}2{,}083$; 41.7\% gap rate), we report unique payload counts and family-level gap rates that reveal sharp heterogeneity between spawn strategies, directly informing which rule families require improvement.
\item \textbf{Normalization control.} We include a 20-request pilot (10 rows in \texttt{attack.json}, 10 in \texttt{benign.json}) comparing a TF-IDF + logistic-regression surrogate on raw versus RFC-style canonicalized strings, observing identical 42.9\% test accuracy, which cautions against equating single-pass decoding with meaningful policy gap reduction.
\end{enumerate}

\subsection{Paper Organization}
Section~\ref{sec:background} situates mutation-based SQLi, discrepancy attacks, and the defensive evaluation landscape. Section~\ref{sec:methodology} defines metrics, outcome semantics, and audit toolchain stages. Section~\ref{sec:experiments} answers RQ1--RQ3 with corpus statistics, aggregate and family-level coverage gap tables, and the normalization pilot. Section~\ref{sec:discussion} discusses validity, defender implications, ethics, and future work.
